\chapter{Expected Outcome}

The expected outcome of the ContextDrive project is a user-friendly, smartphone-based advanced driver-assistance system (ADAS) that assists drivers in making informed, safe decisions using real-time camera, GPS, and environmental information. The system integrates edge-based computer vision, continuous telemetry monitoring, environmental intelligence, deterministic risk assessment, and transparent explainability into a single mobile application. The vision and object detection module is expected to identify potential hazards, such as leading vehicles, by efficiently processing live camera frames on-device. The risk assessment module is expected to provide practical safety evaluations by analyzing vehicle speed, object proximity, and adverse weather conditions, helping drivers maintain safe following distances and adapt to changing environments. The transparent explainability module is expected to improve trust and interpretability by explicitly identifying the critical factors (e.g., high speed in low visibility) that triggered a safety warning. Overall, the project is expected to provide an integrated smart driving platform that supports data-driven risk evaluation, enhanced situational awareness, and context-aware safety guidance.

The system is also expected to provide a seamless and accessible interface through which drivers can view their live camera feed, real-time risk scores, bounding boxes for detected hazards, and actionable safety recommendations. By bringing these functionalities together in a single smartphone application without the need for expensive dedicated hardware, the project aims to reduce the complexity of driving safety systems and support drivers in adopting safer, more proactive driving habits.